%%%%%%%%%%%%%%%%%%%%%%%%%%%%%%%%%%%%%%%%%
% Jacobs Landscape Poster
% LaTeX Template
% Version 1.1 (14/06/14)
%
% Created by:
% Computational Physics and Biophysics Group, Jacobs University
% https://teamwork.jacobs-university.de:8443/confluence/display/CoPandBiG/LaTeX+Poster
% 
% Further modified by:
% Nathaniel Johnston (nathaniel@njohnston.ca)
%
% This template has been downloaded from:
% http://www.LaTeXTemplates.com
%
% License:
% CC BY-NC-SA 3.0 (http://creativecommons.org/licenses/by-nc-sa/3.0/)
%
%%%%%%%%%%%%%%%%%%%%%%%%%%%%%%%%%%%%%%%%%

%----------------------------------------------------------------------------------------
%	PACKAGES AND OTHER DOCUMENT CONFIGURATIONS
%----------------------------------------------------------------------------------------

\documentclass[final]{beamer}

\usepackage[scale=1.24]{beamerposter} % Use the beamerposter package for laying out the poster


\usetheme{confposter} % Use the confposter theme supplied with this template

\setbeamercolor{block title}{fg=dblue,bg=white} % Colors of the block titles
\setbeamercolor{block body}{fg=black,bg=white} % Colors of the body of blocks
\setbeamercolor{block alerted title}{fg=white,bg=dblue!70} % Colors of the highlighted block titles
\setbeamercolor{block alerted body}{fg=black,bg=dblue!10} % Colors of the body of highlighted blocks
% Many more colors are available for use in beamerthemeconfposter.sty

% Top-right contact links in the header
\newcommand{\contactlinks}{%
  \href{mailto:nnepal2@ncsu.edu}{nnepal2@ncsu.edu}
}
\addtobeamertemplate{headline}{}{%
  \begin{tikzpicture}[remember picture,overlay][scale=1.3]
    \node[anchor=north east, inner sep=6pt, fill=white, rounded corners=2pt, draw=white]
      at ([xshift=-10mm,yshift=-115mm]current page.north east)
      {\small \contactlinks};
  \end{tikzpicture}%
}
%-----------------------------------------------------------
% Define the column widths and overall poster size
% To set effective sepwid, onecolwid and twocolwid values, first choose how many columns you want and how much separation you want between columns
% In this template, the separation width chosen is 0.024 of the paper width and a 4-column layout
% onecolwid should therefore be (1-(# of columns+1)*sepwid)/# of columns e.g. (1-(4+1)*0.024)/4 = 0.22
% Set twocolwid to be (2*onecolwid)+sepwid = 0.464
% Set threecolwid to be (3*onecolwid)+2*sepwid = 0.708

\newlength{\sepwid}
\newlength{\onecolwid}
\newlength{\twocolwid}
\newlength{\threecolwid}
\setlength{\paperwidth}{48in} % A0 width: 46.8in
\setlength{\paperheight}{36in} % A0 height: 33.1in
\setlength{\sepwid}{0.003\paperwidth} % Separation width (white space) between columns
\setlength{\onecolwid}{0.31\paperwidth} % Width of one column
\setlength{\twocolwid}{0.28\paperwidth} % Width of two columns
\setlength{\threecolwid}{0.31\paperwidth} % Width of three columns
\setlength{\topmargin}{-0.9in} % Reduce the top margin size
%-----------------------------------------------------------

\setbeamerfont{block alerted title}{size=\small}
\setbeamerfont{block alerted body}{size=\small}
\usepackage{graphicx}  % Required for including images

% Use scalable fonts to avoid cmex size substitution warnings
\usepackage{lmodern}
\usepackage{exscale}
\usefonttheme{professionalfonts}

\usepackage{booktabs} % Top and bottom rules for tables

\usepackage{tikz}
\usepackage{tikz-cd}

\usepackage{biblatex}
\addbibresource{biblio.bib}

% Show all authors in bibliography (no et al.) with first name as initial
\ExecuteBibliographyOptions{maxbibnames=99, maxcitenames=99, giveninits=true}

%\usepackage[utf8]{inputenc}

\usetikzlibrary{matrix,shapes,arrows,positioning,chains,decorations.pathmorphing,snakes,calc}

\DeclareMathOperator{\rk}{rk}
\DeclareMathOperator{\cusp}{cusp}
\DeclareMathOperator{\Rel}{Rel}
\DeclareMathOperator{\Poin}{Poin}
\DeclareMathOperator{\Tor}{Tor}
\DeclareMathOperator{\Ext}{Ext}

\newcommand{\M}{\mathsf{M}}
\newcommand{\FY}{\mathsf{FY}}
\newcommand{\grp}{\mathsf{G}}
\newcommand{\stab}{\mathsf{Stab}}
%\newcommand{\B}{\mathsf{B}} let us only write \U_{n,n} for boolean matroids
\newcommand{\N}{\mathsf{N}}
\newcommand{\U}{\mathsf{U}}
\newcommand{\cI}{\mathcal{I}}
\newcommand{\symm}{\mathfrak{S}}
\newcommand{\A}{\mathrm{A}}
\newcommand{\Q}{\mathbb{Q}}
\newcommand{\R}{\mathbb{R}}
\newcommand{\Z}{\mathbb{Z}}

\newcommand{\LL}{\mathsf{\Lambda}}
\newcommand{\Hilb}{\operatorname{Hilb}}
\newcommand{\rank}{\operatorname{rk}}
\newcommand{\corank}{\operatorname{crk}}
\newcommand{\cl}{\operatorname{cl}}
\renewcommand{\H}{\mathrm{H}}
\newcommand{\CH}{\mathrm{CH}}
\newcommand{\aug}{\operatorname{aug}}
\newcommand{\IH}{\mathrm{IH}}
\newcommand{\uH}{\underline{\mathrm{H}}}
\newcommand{\uCH}{\underline{\mathrm{CH}}}
\newcommand\size[1]{|#1|}
\DeclareMathOperator{\Pal}{\mathrm{Pal}}
\DeclareMathOperator{\Poly}{\mathrm{Poly}}

\newcommand{\cL}{\mathcal{L}}
\newcommand{\Rep}{\operatorname{Rep}}
\newcommand{\gr}{\operatorname{gr}}
\newcommand{\VRep}{\operatorname{VRep}}
\newcommand{\Ind}{\operatorname{Ind}}
\newcommand{\Res}{\operatorname{Res}}
\newcommand{\Aut}{\operatorname{Aut}}
\newcommand{\col}{\operatorname{col}}
\newcommand{\asc}{\operatorname{asc}}
\newcommand{\bad}{\operatorname{bad}}
\newcommand{\des}{\operatorname{des}}

\newcommand{\Ver}{7}
\newcommand{\Vera}{2}
\newcommand{\Hor}{15}
\newcommand{\Hora}{20}
\newcommand{\Horb}{3}
\newcommand{\Horc}{4.5}
\newcommand{\Hord}{6.3}
\newcommand{\Hore}{11}
\newcommand{\Horf}{7.7}
\newcommand{\Tone}{1.85}

\newcommand{\CC}{\mathbb{C}}
\newcommand{\TT}{\mathbb{T}}

\newcommand{\fg}{\mathfrak{g}}

\newcommand{\cF}{\mathcal{F}}
\newcommand{\cU}{\mathcal{U}}
\newcommand{\cO}{\mathcal{O}}

\newcommand{\bw}{\mathbf{w}}
\newcommand{\bP}{\mathbf{P}}
\newcommand{\be}{\mathbf{e}}

\newcommand{\sT}{\mathsf{T}}

\newcommand{\Db}{D^{\mathrm{b}}}
\newcommand{\GL}{\mathbf{GL}}
\newcommand{\IC}{\mathrm{IC}}



\newcounter{sarrow}
\newcommand\xrsquigarrow[1]{%
\stepcounter{sarrow}%
\begin{tikzpicture}[decoration=snake]
\node (\thesarrow) {\strut#1};
\draw[->,decorate] (\thesarrow.south west) -- (\thesarrow.south east);
\end{tikzpicture}%
}

\tikzset{
    ncbar angle/.initial=90,
    ncbar/.style={
        to path=(\tikztostart)
        -- ($(\tikztostart)!#1!\pgfkeysvalueof{/tikz/ncbar angle}:(\tikztotarget)$)
        -- ($(\tikztotarget)!($(\tikztostart)!#1!\pgfkeysvalueof{/tikz/ncbar angle}:(\tikztotarget)$)!\pgfkeysvalueof{/tikz/ncbar angle}:(\tikztostart)$)
        -- (\tikztotarget)
    },
    ncbar/.default=0.5cm,
}

\tikzset{square left brace/.style={ncbar=0.5cm}}
\tikzset{square right brace/.style={ncbar=-0.5cm}}

\tikzset{round left paren/.style={ncbar=0.5cm,out=120,in=-120}}
\tikzset{round right paren/.style={ncbar=0.5cm,out=60,in=-60}}

%----------------------------------------------------------------------------------------
%	TITLE SECTION 
%----------------------------------------------------------------------------------------

\title{A Deletion Formula for Inverse Kazhdan--Lusztig Polynomials of Matroids} % Poster title

\author{Tom Braden, Luis Ferroni, Jacob P. Matherne, and \color{nred}{Nutan Nepal}} % Author(s)

\institute{Department of Mathematics, North Carolina State University} % Institution(s)

%----------------------------------------------------------------------------------------

\begin{document}

\addtobeamertemplate{block end}{}{\vspace*{2ex}} % White space under blocks
\addtobeamertemplate{block alerted end}{}{\vspace*{2ex}} % White space under highlighted (alert) blocks

\setlength{\belowcaptionskip}{2ex} % White space under figures
\setlength\belowdisplayshortskip{2ex} % White space under equations

\begin{frame}[t] % The whole poster is enclosed in one beamer frame

\begin{columns}[t] % The whole poster consists of three major columns, the second of which is split into two columns twice - the [t] option aligns each column's content to the top


\begin{column}{\sepwid}\end{column} % Empty spacer column

\begin{column}{\onecolwid} % The first column
\vspace{-17mm}
\begin{block}{Polynomial invariants of matroids}
    To a loopless matroid $\M$, there
    are several important associated polynomial invariants:
    \begin{enumerate}[1]
        \item The characteristic polynomial \(\chi_\M(x)\),
        \item The Chow polynomial \(\uH_\M(x)\) \cite{ferroni-matherne-stevens-vecchi},
        \item The Kazhdan--Lusztig polynomial
        \(\mathrm{P}_\M(x)\) \cite{EPW16},
        \item The inverse Kazhdan--Lusztig polynomial
        \(\mathrm{Q}_\M(x)\) \cite{gao-xie}.
    \end{enumerate}
    These polynomials have the following algebraic interpretations:
    \begin{enumerate}[1']
        \item \(\chi_\M(x)\) encodes the Poincaré polynomial of the
        Orlik--Solomon algebra of \(\M\),
        \item \(\uH_\M(x)\) is the Poincaré polynomial of the Chow ring of \(\M\)
        \cite{feichtner-yuzvinsky, ferroni-matherne-stevens-vecchi},
        \item \(\mathrm{P}_\M(x)\) is the Poincaré polynomial of the local
            intersection cohomology module of $\M$
            \cite{braden-huh-matherne-proudfoot-wang},
        \item The coefficient of $x^k$ in $\mathrm{Q}_{\M}(x)$ is the
        multiplicity of the trivial module in the degree
        $\rk(\M) - 2k$ piece of the Rouquier complex of
        $\M$ \cite{braden-huh-matherne-proudfoot-wang}.
    \end{enumerate}
\end{block}
\vspace{-15mm}
\begin{alertblock}{}
    \emph{Definition.} The inverse Kazhdan--Lusztig polynomial $\mathrm{Q}_\M(x)$ 
    of a loopless matroid $\M$
    is the unique polynomial satisfying
    \[
        \sum_{F \in \cL(\M)}
        (-1)^{\rk(F)}\mathrm{Q}_{\M|_F}(x)
        \cdot \mathrm{P}_{\M/F}(x) = \delta_{\M}
    \]
    which is 1 if $\M$ is the empty matroid and 0 otherwise.
\end{alertblock}
\vspace{-15mm}
\begin{block}{Examples}
    Let $\M$ be the uniform matroid $U_{3,4}$, the
    graphic matroid on the cycle $C_4$.
    \vspace{-20mm}
    \begin{columns}[t]
        \column{0.28\textwidth}
        \begin{center}
        \begin{tikzpicture}[scale=6]
            % vertices as filled dots
            \fill (0,0) circle[radius=1.2pt];
            \fill (1,0) circle[radius=1.2pt];
            \fill (1,1) circle[radius=1.2pt];
            \fill (0,1) circle[radius=1.2pt];
            % edges of C4
            \draw[red, very thick] (0,0) -- (1,0) node[midway, below]{\(e\)};
            \draw (1,0) -- (1,1) -- (0,1) -- (0,0);
        \end{tikzpicture}
        \end{center}
        \column{0.68\textwidth}
        \begin{enumerate}
            \item $\chi_\M(x) = x^3 - 4x^2 + 6x - 3$.
            \item $\mathrm{P}_\M(x) = 2x + 1$.
            \item $\mathrm{Q}_\M(x) = 2x + 3$.
        \end{enumerate}
    \end{columns}
\end{block}
\vspace{-10mm}
After contracting the edge \(e\), we get the matroid \(\M/e\) on \(C_3\):
\vspace{-10mm}
\begin{columns}[t]
    \column{0.28\textwidth}
    \begin{center}
    \begin{tikzpicture}[scale=6.5]
        % vertices as filled dots
        \fill (-0.7,1) circle[radius=1.2pt];
        \fill (0.7,1) circle[radius=1.2pt];
        \fill (0,0) circle[radius=1.2pt];
        \draw (-0.7,1) -- (0.7,1) -- (0,0) -- cycle;
    \end{tikzpicture}
    \end{center}
    \column{0.68\textwidth}
    \begin{enumerate}
        \item $\chi_{\M/e}(x) = x^2 - 3x + 2$.
        \item $\mathrm{P}_{\M/e}(x) = 1$.
        \item $\mathrm{Q}_{\M/e}(x) = 2$.
    \end{enumerate}
\end{columns}
\vspace{5mm}
After deleting the edge \(e\), we get the matroid
\(\M\setminus e\) on the tree:
\vspace{-10mm}
\begin{columns}[t]
    \column{0.28\textwidth}
    \begin{center}
    \begin{tikzpicture}[scale=5]
        \fill (0,0) circle[radius=1.2pt];
        \fill (1,0) circle[radius=1.2pt];
        \fill (1,1) circle[radius=1.2pt];
        \fill (0,1) circle[radius=1.2pt];
        % edges of C3
        \draw (1,0) -- (1,1) -- (0,1) -- (0,0);
    \end{tikzpicture}
    \end{center}
    \column{0.68\textwidth}
    \begin{enumerate}
        \item $\chi_{\M\setminus e}(x) = x^3 - 3x^2 + 3x - 1$.
        \item $\mathrm{P}_{\M\setminus e}(x) = 1$.
        \item $\mathrm{Q}_{\M\setminus e}(x) = 1$.
    \end{enumerate}
\end{columns}
\end{column} % End of the first column






\begin{column}{\twocolwid} % The second col

\vspace{-17mm}
\begin{block}{Deletion formulas for the invariants}
The aforementioned invariants satisfy the following formulas:
\vspace{-10mm}
\begin{enumerate}
    \item $\chi_\M(x) = \chi_{\M\setminus e}(x) - \chi_{\M/e}(x)$.
    \item $\displaystyle\mathrm{P}_\M(x) = \mathrm{P}_{\M\setminus e}(x) - x \mathrm{P}_{\M/e}(x) +
    \sum_{F \in S}{\tau(\M/{F\cup e})x^{\corank(F)/2}\cdot
    \mathrm{P}_{\M|_F}(x)}$
    where $S := \{F \in \cL(\M) \mid e \notin F \text{ and } F \cup e \in \cL(\M)\}$
    and $\tau(\M)$ is the coefficient of $x^{\lfloor\rank(\M)/2\rfloor}$ in the
    \(\mathrm{P}_\M(x)\) \cite{braden-vysogorets}.
\end{enumerate}
\begingroup
\setbeamercolor{block alerted title}{fg=white,bg=nred!70}
\setbeamercolor{block alerted body}{fg=black,bg=nred!10}
\begin{alertblock}{}
    \emph{Main theorem \cite{braden2025inversekazhdanlusztigpolynomialsmatroids}.} 
    The inverse Kazhdan--Lusztig polynomial satisfies the following deletion formula:
    \begin{align*}
    \mathrm{Q}_\M(x)
    & =
    \mathrm{Q}_{\M\setminus e}(x) + (1 + x) \mathrm{Q}_{\M/e}(x) - \\
    &\hspace{5em}\sum_{F \in T}{\tau(\M|_F/ e)x^{\rank(F)/2}\cdot \mathrm{Q}_{\M/F}(x)}
    \end{align*}
    where $T := \{F \in \cL(\M) \mid e \in F \text{ and } F \setminus e \notin \cL(\M)\}$.
\end{alertblock}
\endgroup
\vspace{-10mm}
    \begin{figure}
    \scalebox{0.4}{
    \begin{tikzpicture}[scale=5.5]

        % Level 1 (topmost node)
        \node at (0, 4) (A) {\{1,2,3,4\}};
        
        % Level 2 (second row of nodes)
        \node at (-5, 2) (B1) {\{1,2\}};
        \node at (-3, 2) (B2) {\{1,3\}};
        \node at (-1, 2) (B3) {\{1,4\}};
        \node at (1, 2) (B4) {\{2,3\}};
        \node at (3, 2) (B5) {\{2,4\}};
        \node at (5, 2) (B6) {\{3,4\}};
        
        % Level 3 (third row of nodes)
        \node at (-3, 0) (C1) {\{1\}};
        \node at (-1, 0) (C2) {\{2\}};
        \node at (1, 0) (C3) {\{3\}};
        \node at (3, 0) (C4) {\{4\}};
        
        % Level 4 (bottommost node)
        \node at (0, -2) (D) {$\varnothing$};
        
        % Arrows from top to second row
        \draw[->] (A) -- (B1);
        \draw[->] (A) -- (B2);
        \draw[->] (A) -- (B3);
        \draw[->] (A) -- (B4);
        \draw[->] (A) -- (B5);
        \draw[->] (A) -- (B6);
        
        % Arrows from second to third row
        \draw[->] (B1) -- (C1);
        \draw[->] (B1) -- (C2);
        
        \draw[->] (B2) -- (C1);
        \draw[->] (B2) -- (C3);
        
        \draw[->] (B3) -- (C1);
        \draw[->] (B3) -- (C4);
        
        \draw[->] (B4) -- (C2);
        \draw[->] (B4) -- (C3);
        
        \draw[->] (B5) -- (C2);
        \draw[->] (B5) -- (C4);
        
        \draw[->] (B6) -- (C3);
        \draw[->] (B6) -- (C4);
        
        % Arrows from third row to bottom
        \draw[->] (C1) -- (D);
        \draw[->] (C2) -- (D);
        \draw[->] (C3) -- (D);
        \draw[->] (C4) -- (D);
        
        \end{tikzpicture}
    }
\end{figure}
\vspace{-5mm}
From the lattice of flats of $U_{3,4}$ above,
we can see that the set \(T\) consists of just the top flat for which
\(\tau(\M|_F/ e) = 0\) and the
deletion formula reduces to
\[
\mathrm{Q}_\M(x) = 1 + (1 +x)\cdot 2 - 0 = 2x + 3.
\]
\end{block}

\vspace*{-17mm}

\begin{block}{Application: Counterexample to a log-concavity conjecture}
A polynomial \(\sum_{i=0}^d{a_i x^i} \in \mathbb{R}_{\geq 0}[x]\)
is log-concave if \(a_k^2 \geq a_{k-1} a_{k+1}\) for all \(0<k<d\).
The log-concavity of the polynomial invariants of matroids
is a topic of ongoing research.
\begin{enumerate}
    \item The log-concavity of \(\chi_\M(x)\)
    (Heron--Rota--Welsh conjecture) was proved by Adiprasito, Huh, and Katz in
    \cite{adiprasito-huh-katz}.
    \item \textbf{Conjecture:} The polynomials
    \(\mathrm{P}_\M(x)\) and \(\mathrm{Q}_\M(x)\)
        are log-concave \cite{gedeon-proudfoot-young-survey,gao-xie}.
\end{enumerate}
In their proof of the log-concavity of the inverse Kazhdan--Lusztig polynomial
of paving matroids, Xie and Zhang conjectured that the
Hadamard product of the inverse Kazhdan--Lusztig polynomial of degree $d$
with the polynomial \((1+x)^d\) is real-rooted.
This conjecture, if true,
would imply the log-concavity of the inverse Kazhdan--Lusztig polynomial
of all matroids. 
\end{block}
\end{column} % End of the second column

\begin{column}{\threecolwid} % The third col

\vspace{-20mm}
\begin{block}{}
We use the deletion formula to provide a counterexample
to this conjecture.
Let $\M$ be the rank 2 matroid on the ground set $\{1,2,\ldots, 21\}$.
\begin{figure}
    \scalebox{0.4}{
    \begin{tikzpicture}[scale=6]
    
            % Level 1 (topmost node)
            \node at (0, 4) (A) {\{1,2,\ldots, 21\}};
            
            % Level 2 (second row of nodes)
            \node at (-5, 2) (B1) {\{1,2,3,4\}};
            \node at (-3, 2) (B2) {\{5,6,7,8\}};
            \node at (-1, 2) (B3) {\{9,10,11,12\}};
            \node at (1, 2) (B4) {\{13,14,15\}};
            \node at (3, 2) (B5) {\{16,17,18\}};
            \node at (5, 2) (B6) {\{19,20,21\}};
            
            % Level 3 (third row of nodes)
            \node at (0, 0) (C) {\(\varnothing\)};
            
            % Arrows from second to third row
            \draw[->] (B1) -- (C);
            \draw[->] (B2) -- (C);
            \draw[->] (B3) -- (C);
            \draw[->] (B4) -- (C);
            \draw[->] (B5) -- (C);
            \draw[->] (B6) -- (C);

            % Arrows from top to second row
            \draw[->] (A) -- (B1);
            \draw[->] (A) -- (B2);
            \draw[->] (A) -- (B3);
            \draw[->] (A) -- (B4);
            \draw[->] (A) -- (B5);
            \draw[->] (A) -- (B6);

            
            \end{tikzpicture}
    }
\end{figure}
\vspace{-10mm}
Using the deletion formula,
we can compute the inverse Kazhdan--Lusztig polynomial of the
dual matroid $\M^*$. We find the Hadamard product
\(\mathrm{Q}_{\M^*}(x) \star (1+x)^{\text{deg}(\mathrm{Q}_{\M^*}(x))}\)
to be
\begin{align*}
71162 x^9 + 2093958 x^8 + 12614544 x^7 + 27186264 x^6 +
27111294 x^5 + \\
13439034 x^4 + 3294228 x^3 + 371628 x^2 +
16110 x + 163
\end{align*}
which we can see, by using SageMath, is not real-rooted.
Note that this does not contradict the conjecture regarding
the log-concavity of the polynomials
\(\mathrm{P}_\M(x)\) and \(\mathrm{Q}_\M(x)\).

\end{block}
\vspace{-10mm}
\begin{alertblock}{References}

\printbibliography
\end{alertblock}
\end{column} % End of the third column

\end{columns} % End of all the columns in the poster

\end{frame} % End of the enclosing frame

\end{document}
